\scanpage{86}
\[
\tag{41}f_{l_i}(u)\qquad i\in I,\ u\in\Gamma
\]
where for every $i\in I$, $l_i$ is a path in $\Pi_1(X,x_i)$ from $x_0$ to $x_i$, i.e. an element of the ``discrete'' part of the torsor (41).

On the other hand, the action of $\Gamma$ on $\widehat\Pi_1(X,x_0)$ is known when one \unsure{knows|knows} its action on generators $\lambda_j$ \inserted{($j\in J$)} of $\Pi_1(X,x_0)$, ie. if one \unsure{knows|knows} the
\[
\tag{43}f_{\lambda_j}(u)\quad\text{or}\quad g_{\lambda_j}(u)\qquad j\in J,\ u\in\Gamma.
\]
\deleted{But this} The data (42)(43) therefore suffice to determine the action of $\Gamma$ on $\Pi_1(X,I)$ --- but this description, \deleted{at the} \deleted{dis} if it has the advantage of an economy, at least in appearance, has the disadvantage of referring to arbitrary choices, the worst being the choice \unsure{of a|of a} \deleted{origin} \inserted{base point} $x_0$ \deleted{to which} \inserted{to which} one is supposed to reduce everything --- which seems artificial, and in the long run I think impracticable, \deleted{\unsure{and|is}} and most particularly in situations involving a group of \unsure{symmetry|symmetries} $\mathcal G$ as above, when one has taken $I$ stable under $\mathcal G$ ---
\margin{which must be \unsure{expressed} \unsure{linearly|literally}, \unsure{resp.|resp} principle, \unsure{described on page 8!|said on page 8'}}
\margin{85}
